\scanpage{91}

\[
(28)\quad
\left\{\begin{aligned}
\widetilde{\Pi}_{0,3}=\Pi_{0,3}\cdot\{1,\tau\}&\longrightarrow\widetilde{\Pi}'_{0,3}\\
\tau_0,\tau_1,\tau_\infty&\longmapsto
\sigma'_0,\sigma'_1,\sigma'_\infty
\end{aligned}\right.
\]

Ceci posé, on définit un \uline{iso.} \deleted{\illegible{1}}
\[
(29)\quad
\begin{aligned}
\widetilde{\Pi}'_{0,3}&\xrightarrow{\sim}\mathfrak S_{0,3}\\
(\sigma'_0,\sigma'_1,\sigma'_\infty)
&\longrightarrow
\left(\deleted{\unsure{\sigma^{\prime}_1}}
\underset{\substack{\inserted{\text{sur }\dot\tau'_\infty}\\\Vert\\\inserted{\dot\sigma_\infty\tau}}}{\tau'_\infty},
\underset{\substack{\inserted{\text{sur }\dot\tau'_0}\\\Vert\\\inserted{\dot\sigma_0\tau}}}{\tau'_0},
\underset{\inserted{\text{sur }\tau}}{\tau_\infty}\right)
\end{aligned}
\]
notés aussi $r_0,r_1,r_\infty$ (\uline{réflexions} fondamentales)

\margin{Cet iso. est associé à un homomorphisme canonique
$\varphi:\widetilde{\Pi}_{0,3}=\Pi_{0,3}\cdot\{1,\tau\}\xrightarrow{\text{épi}}\mathfrak S_{0,3}$
induisant un hom. can.
$\varphi:\Pi_{0,3}\xrightarrow{\text{épi}}\mathfrak S_{0,3}^{+}$
permettant d'identifier $\mathfrak S_{0,3}$ et $\mathfrak S_{0,3}^{+}$
aux quotients des pr.~\unsure{mbres.|membres} par les relations
$\ell_1^2=\ell_\infty^3=1$ \ldots}

qui induit un iso.
\[
(30)\quad
\left\{\begin{aligned}
\Pi'_{0,3}&\xrightarrow{\sim}\deleted{\unsure{\Pi}}\mathfrak S_{0,3}^{+}\\
\rho_s=\ell'_0&\longmapsto\varepsilon_0=\sigma_\infty\rho\\
\sigma=\ell'_1&\longmapsto\sigma_\infty\\
\rho_f=\ell'_\infty&\longmapsto\rho^{-1}
\end{aligned}\right.
\]

\margin{\uline{NB}
$\left\{\begin{aligned}
\tau_\infty\tau'_0&=\varepsilon_0=\ell'_0\\
\tau'_\infty\tau_\infty&=\sigma_\infty=\ell'_1\\
\tau'_0\tau'_\infty&=\rho^{-1}=\ell'_\infty
\end{aligned}\right.$
i.e.
$\begin{aligned}
r_\infty r_1&=\ell'_0\;(=\varepsilon_0)\\
r_0r_\infty&=\ell'_1\;(=\sigma_\infty)\\
r_1r_0&=\ell'_\infty\;(=\rho^{-1})
\end{aligned}$}

On trouve donc que $\mathfrak S_{0,3}$ a des \deleted{\unsure{gén}} relations génératrices
\[
(31)\quad
{\tau'_\infty}^{2}={\tau'_0}^{2}=\tau_\infty^2
=\bigl(\underbrace{\tau'_\infty\tau_\infty}_{\sigma_\infty}\bigr)^2
=\bigl(\underbrace{\tau'_0\tau'_\infty}_{\rho^{-1}}\bigr)^3=1,
\]
et $\mathfrak S_{0,3}^{+}$ a des relations génératrices
\[
(32)\quad \ell'_\infty \ell'_1\ell'_0={\ell'_1}^{2}={\ell'_\infty}^{3}=1
\]
ou aussi simplement, en éliminant $\ell'_0$ et posant
$\ell'_1=\sigma_\infty$ et ${\ell'_\infty}^{-1}=\rho$
\[
(33)\quad\boxed{\sigma_\infty^2=\rho^3=1}
\]

Pour mémoire, on va exprimer les éléments remarquables de $\mathfrak S_{0,3}^{+}$ en termes de $\sigma_\infty,\rho$ :

\scanpage{92}
\mspage{335}

\[
(34)\quad\left\{\begin{aligned}
\varepsilon_0&=\sigma_\infty\rho\ \deleted{\illegible{1}}&
\ell_0&=\sigma_\infty\rho\sigma_\infty\rho=(\sigma_\infty\rho)^2\\
\varepsilon_1&=\rho\sigma_\infty&
\ell_1&=\rho\sigma_\infty\rho\sigma_\infty=(\rho\sigma_\infty)^2\\
\varepsilon_\infty&=\rho^{-1}\sigma_\infty\rho^{-1}&
\ell_\infty&=(\rho^{-1}\sigma_\infty\rho^{-1})^2
=\rho^{-1}\sigma_\infty\rho^{-1}\sigma_\infty\rho^{-1}
\end{aligned}\right.,
\]
\[
(35)\quad\left\{\begin{aligned}
\sigma_0&=\rho^{\deleted{\unsure{2}}}\sigma_\infty\rho^{\deleted{\unsure{2}}-1}
\quad(=\deleted{\unsure{\rho}}\varepsilon_1\rho^{-1}=\rho\varepsilon_1^{-1})\\
\sigma_1&=\rho^{-1}\sigma_\infty\rho^{\deleted{\unsure{2}}}
\quad(=\rho^{-1}\varepsilon_0\deleted{\unsure{\rho}}=\varepsilon_\infty\rho^{-1})
\end{aligned}\right.
\]

et dans $\mathfrak S_{0,3}^{-}$ on a, en fonction de $\tau'_\infty,\tau'_0,\tau_\infty$,
les expressions de \deleted{$\unsure{\tau_0,\tau_1,\tau_\infty,\tau'_0,\tau'_1,\tau'_\infty}$}
\inserted{$\tau_0,\tau_1,(\tau_\infty,\tau'_0),\tau'_1,(\tau'_\infty)$}
\[
(36)\quad\left\{\begin{aligned}
\tau_0&=\deleted{\unsure{\tau_0\tau'_0\tau_\infty\tau'_\infty}}\rho\tau_\infty\rho^{-1}
=\tau'_\infty\tau'_0\tau_\infty\tau'_0\tau'_\infty
\quad(=r_0r_1r_\infty r_1r_0)\\
\tau_1&=\rho^{-1}\tau_\infty\rho
=\tau'_0\tau'_\infty\tau_\infty\tau'_\infty\tau'_0
\quad(=r_1r_0r_\infty r_0r_1)\\
\tau'_1&=\rho\tau'_0\rho^{-1}
=\tau'_\infty\deleted{\unsure{\tau'_0}}\tau'_0\tau'_\infty
\quad(=r_0r_1r_0)
\end{aligned}\right.
\]
et
\[
\left\{\begin{aligned}
\tau'_\infty&=r_0\\
\tau'_0&=r_1\\
\tau_\infty&=r_\infty.
\end{aligned}\right.
\]

\uline{IV}) Relation avec
$\mathrm{Gl}(2,\mathbb Z)'=\mathrm{Gl}(2,\mathbb Z)/\mathord{\pm}1
\simeq\mathrm{GP}(1,\mathbb Z)$.

Ce groupe opère sur le demi-plan de Poincaré
\[
D=\{z\in\mathbb C\mid\Im z>0\}
\]
par
\[
(37)\quad\left\{\begin{aligned}
u_g(z)&=\frac{az+b}{cz+d}&\text{si }ad-bc&=+1\\
u_g(z)&=\frac{a\bar z+b}{c\bar z+d}&\text{si }ad-bc&=-1
\end{aligned}\right.
\qquad
\begin{gathered}
g=\begin{pmatrix}a&b\\c&d\end{pmatrix}\\
a,b,c,d\in\mathbb Z\\
ad-bc\in\{\pm1\}
\end{gathered}
\]

On a dans $D$ un domaine fondamental hachuré, qui correspond au triangle
sphérique $(-\bar j,0,\tfrac12)$ par l'\unsure{homomorphisme|homéomorphisme}
de \unsure{revêt.} naturel, $i$ correspondant à $\tfrac12$,
$-\bar j$ à $-\bar j$, et le point à l'infini à $0$.
\deleted{Ainsi \unsure{on trouve}}\inserted{Les réflexions}
par rapport aux trois côtés de ce triangle sont donnés par
\[
(38)\quad
\deleted{\unsure{\tau_0}}\unsure{r_0|\tau_0}(z)=\frac1{\bar z},\qquad
\unsure{r_1|\tau_1}(z)=1-\bar z,\qquad
\unsure{r_\infty|\tau_\infty}(z)=-\bar z
\]
\figure{\begin{tikzpicture}[x=2.2cm,y=2.2cm,line cap=round,line join=round,font=\scriptsize]
\draw[line width=.4pt] (-1.55,0) .. controls (-.7,.04) and (.5,-.025) .. (1.85,.01);
\draw[line width=.4pt] (-1,0) arc[start angle=180,end angle=0,radius=1];
\draw[line width=.5pt] (0,-.48)--(0,2.85);
\draw[line width=.5pt] (.5,-.65)--(.5,2.78);
\begin{scope}
\clip (0,1) arc[start angle=90,end angle=60,radius=1]--(.5,2.78)--(0,2.85)--cycle;
\foreach \h in {.96,1.24,1.52,1.80,2.08,2.36,2.64}
  \draw[line width=.35pt] (0,\h)--(.5,{\h+.30});
\end{scope}
\draw[line width=.7pt] (0,1)--(.075,1)--(.075,1.075)--(0,1.075);
\fill (.5,.8660254) circle[radius=.025];
\node[inner sep=0pt] at (-1.04,-.14) {$-1$};
\node[inner sep=0pt] at (-.10,-.14) {$0$};
\node[inner sep=0pt] at (.43,-.18) {$\frac12$};
\node[inner sep=0pt] at (1.04,-.14) {$1$};
\node[inner sep=0pt] at (-.13,1.12) {$i$};
\node[inner sep=0pt] at (.64,.94) {$-\bar j$};
\node[inner sep=0pt,anchor=east,rotate=16] at (-.07,1.94) {$\dot r_\infty=\dot\tau_\infty$};
\node[inner sep=0pt,anchor=west] at (.56,2.25) {$\dot\tau'_0=\dot r_1$};
\node[inner sep=0pt,rotate=-63] at (.23,.63) {$\dot\tau'_\infty=\dot r_0$};
\end{tikzpicture}}

\deleted{\unsure{Ils déterminent}} et les produits deux à deux de ces réflexions

\scanpage{93}

sont
\[
(39)\quad\left\{\begin{aligned}
\unsure{r_\infty r_1|\tau_\infty\tau_1}(z)&=\ell'_0(z)=z-1\;\deleted{\unsure{;}}\\\unsure{r_0r_\infty|\tau_0\tau_\infty}(z)&=\ell'_1(z)=-\frac1z\\
\unsure{r_1r_0|\tau_1\tau_0}(z)&=\ell'_\infty(z)=1-\frac1z=\frac{z-1}{z}
\end{aligned}\right.
\]

On trouve ainsi un isomorphisme
\[
(40)\quad\left\{\begin{aligned}
\mathfrak S_{0,3}&\xrightarrow{\sim}\mathrm{Gl}(2,\mathbb Z)'\simeq
\mathrm{Gl}(2,\mathbb Z)/\mathord{\pm}1\simeq\mathrm{GP}(1,\mathbb Z)\\
r_0=\tau'_\infty&\longmapsto r_0=\begin{pmatrix}0&1\\1&0\end{pmatrix}\\
r_1=\tau'_0&\longmapsto r_1=\begin{pmatrix}-1&1\\0&1\end{pmatrix}\\
r_\infty=\tau_\infty&\longmapsto r_\infty=\begin{pmatrix}1&0\\0&-1\end{pmatrix}
\end{aligned}\right.
\]
\margin{il vaut mieux changer \unsure{des|de} signes \unsure{cf. Formulaire|cf formulaire}}

et une présentation de $\mathrm{GP}(1,\mathbb Z)$ par
\[
(41)\quad r_0^2=r_1^2=r_\infty^2=(r_0r_\infty)^2=(r_1r_0)^3=1.
\]
Cet iso. induit un iso
\[
(42)\quad\left\{\begin{aligned}
\mathfrak S_{0,3}^{+}&\xrightarrow{\sim}
\mathrm{Sl}(2,\mathbb Z)'=\mathrm{Sl}(2,\mathbb Z)/\mathord{\pm}1
\simeq\mathrm{GP}(1,\mathbb Z)^{+}\\
\ell'_1=\sigma_\infty&\longmapsto\begin{pmatrix}0&-1\\1&0\end{pmatrix}\\
\ell'_\infty=\rho^{-1}&\longmapsto\begin{pmatrix}1&-1\\1&0\end{pmatrix}
\end{aligned}\right.
\]
\[
(\text{d'où}\quad\rho\longmapsto\begin{pmatrix}0&1\\-1&1\end{pmatrix}).
\]

On identifie par la suite les éléments de $\mathfrak S_{0,3}$ avec leurs images dans
$\unsure{\mathrm{Gl}|\mathrm{Sl}}(2,\mathbb Z)'\simeq\mathrm{GP}(1,\mathbb Z)$.
Notons le formulaire suivant \unsure{(où on utilise} que pour une matrice
$u=\begin{pmatrix}a&b\\c&d\end{pmatrix}\in\mathrm{Sl}(2,\mathbb Z)$,
\deleted{\illegible{1}}
\deleted{\unsure{pour} $\det u\in\mathbb Z^*=\{\pm1\}$ donc $(\det u)^2=1$}, on a

\scanpage{94}
\mspage{336}

\[
(43)\quad u^{-1}=\begin{pmatrix}d&-b\\-c&a\end{pmatrix}\quad):
\]

\margin{NB. \unsure{2.} \deleted{\unsure{Si det}} $u\in\mathrm{Gl}(2,A)$,
$A$ anneau \unsure{quelc.}, on a
$\det(u)\cdot u^{-1}=\begin{pmatrix}d&-b\\-c&a\end{pmatrix}$}

\deleted{\unsure{Et}}.
\[
(44)\quad\left\{\begin{aligned}
\varepsilon_0=\sigma_\infty\rho&=\begin{pmatrix}1&-1\\0&1\end{pmatrix},&
\varepsilon_0^{-1}&=\begin{pmatrix}1&1\\0&1\end{pmatrix}\\
\varepsilon_1=\rho\sigma_\infty&=\begin{pmatrix}1&0\\1&1\end{pmatrix},&
\varepsilon_1^{-1}&=\begin{pmatrix}1&0\\-1&1\end{pmatrix}\\
\varepsilon_\infty=\rho^{-1}\sigma_\infty\rho^{-1}&=\begin{pmatrix}2&-1\\1&0\end{pmatrix},&
\varepsilon_\infty^{-1}&=\begin{pmatrix}0&1\\-1&2\end{pmatrix}\\
\ell_0=\varepsilon_0^2&=\begin{pmatrix}1&-2\\0&1\end{pmatrix}&
\sigma_0&=\begin{pmatrix}1&-1\\2&-1\end{pmatrix}\\
\ell_1=\varepsilon_1^2&=\begin{pmatrix}1&0\\2&1\end{pmatrix}&
\sigma_1&=\begin{pmatrix}1&-2\\1&-1\end{pmatrix}\\
\ell_\infty=\varepsilon_\infty^2&=\begin{pmatrix}3&-2\\2&-1\end{pmatrix}&
\sigma_\infty&=\begin{pmatrix}0&-1\\1&0\end{pmatrix}
\end{aligned}\right.
\]
\margin{Cf. ``formulaire'' pour liste des relations et d'écritures plus \unsure{complète}.}
\[
(45)\quad\rho=\begin{pmatrix}0&1\\-1&1\end{pmatrix},\qquad
\unsure{\rho^{-1}|\rho^{-1\prime}}=\begin{pmatrix}1&-1\\1&0\end{pmatrix},
\]
\[
(46)\quad \ell_\infty \ell_1\ell_0=\begin{pmatrix}-1&0\\0&-1\end{pmatrix}
\]
l'élément non trivial du centre de $\mathrm{Sl}(2,\mathbb Z)$ qui
(heureusement!) donne $1$ dans $\mathrm{Sl}(2,\mathbb Z)'$

Le groupe $\Pi_{0,3}$ se récupère comme le noyau de l'homomorphisme
\diagram{\begin{tikzcd}[ampersand replacement=\&]
\mathrm{Sl}(2,\mathbb Z)' \arrow[r] \& \mathrm{Sl}(2,\mathbb Z/2\mathbb Z) \arrow[d,phantom,"\simeq"{description,rotate=90}]\\
\& \mathfrak S_3
\end{tikzcd}}
la suite exacte
\[
(47)\quad
\diagram{\begin{tikzcd}[ampersand replacement=\&]
1 \arrow[r] \& \Pi_{0,3}\arrow[r,dash] \& \mathfrak S_{0,3}^{+}\arrow[r,dash] \& \mathcal{T}_{0,3}^{+\unsure{\prime}}\arrow[r,dash]\arrow[d,phantom,"\simeq"{description,rotate=90}] \& 1\\
\&\&\&\mathfrak S_3\&
\end{tikzcd}}
\]
se récupère comme
\[
(48)\quad
\diagram{\begin{tikzcd}[ampersand replacement=\&]
1\arrow[r]\&\Pi(2)'\arrow[r]\arrow[d,phantom,"\simeq"{description,rotate=90}] \& \mathrm{Sl}(2,\mathbb Z)'\arrow[r] \& \mathrm{Sl}(2,\mathbb Z/2\mathbb Z)\arrow[r]\arrow[d,phantom,"\simeq"{description,rotate=90}] \& 1\\
\&\Pi_{0,3}^{\unsure{\beta}}\&\&\mathfrak S_3\&
\end{tikzcd}}
\]

\scanpage{95}

Ça nous fourmille dans \unsure{les mains|la main} de \deleted{\unsure{rigoler}}\inserted{travailler}
dans $\mathrm{Gl}(2,\mathbb Z)$ lui-$\overline{\mathrm m}$, extension centrale de
$\mathfrak S_{0,3}$ --- en fait, cette extension s'interprète comme
$\mathcal{T}_{1,1}$ --- mais nous y reviendrons, ultérieurement\ldots

\uline{V}) Résumons les avatars du $\overline{\mathrm m}$ groupe, que nous avons
\unsure{rencontrés|rencontré} jusqu'à présent

1) $\mathfrak S_{0,3}\overset{\text{dfn}}{=}\operatorname{Aut}_{\mathrm{lac}}(\Pi_{0,3})$
avec \unsure{ss-groupe} $\mathfrak S_{0,3}^{+}$ (respectant orientation)

2) $\mathcal{T}'_{0,4}\simeq\operatorname{Autext}_{\mathrm{lac}}(\Pi_{0,4})'$
aut ext. de $\Pi_{0,4}$ respectant la structure à lacets et fixant la classe de lacets
$\unsure{L_4|L_\lambda}$ (et permutant $L_0,L_1,\unsure{L_\infty|L_3}$)

\deleted{avec $\unsure{\mathcal{T}_{0,4}^{\prime+}}$}
\deleted{$\unsure{A'_{0,4}/A_{0,4}^{\prime0}}\simeq\unsure{\operatorname{Auttop}}(\mathbb P^1_{\mathbb C}\setminus\{0,\unsure{1}\})$}
avec $\mathcal{T}_{0,4}^{\prime+}$ (respectant orientation)

3) $A'_{0,4}/A_{0,4}^{\prime0}$ groupe des composantes connexes de
\[
A'_{0,4}\subset\operatorname{Auttop}
\left(\deleted{\mathbb P^1_{\mathbb C}\setminus\unsure{\{}}
\underset{\substack{\Vert\\\mathbb P^1_{\mathbb C}\setminus\{0,1,\infty;\tfrac12\}}}{U_{0,4}^{\mathrm{top}}}\right)
\]
sous-groupe des autom. qui fixent le point à l'infini $a_{0,4}=\tfrac12$

avec $A_{0,4}^{\prime+}/A_{0,4}^{\prime0}$ (respectant orientation

4) $\Pi_1\bigl((U_{0,3}^{\mathrm{an}},\overbrace{\mathfrak S_3\times\unsure{\mathbb Z_2|\mathbb Z/2}}^{\mathcal{T}_{0,3}}),-\bar j\bigr)$
(décrit à la J. Malgoire)
avec le sous-groupe $\Pi_1\bigl((U_{0,3}^{\mathrm{an}},\mathfrak S_3),-\bar j\bigr)$
$(\mathfrak S_3=\mathcal{T}_{0,3}^{+})$

5) $\Pi_1\bigl((\mathbb P^{1\,\mathrm{an}}_{\mathbb C},0\cdot\{0\}+2\{1\}+3\{\infty\}),-\bar j\bigr)$
(à la J. Malgoire)
avec le sous-groupe $\Pi_1\bigl(\mathbb P^{1\,\mathrm{an}}_{\mathbb C},\unsure{0}\{0\}+2\{1\}+3\{\infty\}\bigr)$

6) Groupe cartographique triangulé \deleted{\illegible{1}} $\Gamma_{\infty,3}$
\[
\Gamma_{\infty,3}\unsure{\}}=\{r_0,r_1,r_\infty\mid
r_0^2=r_1^2=r_\infty^2=(r_0r_\infty)^2=(r_1r_0)^3=1\}
\]
(avec sous-groupe cart. triangulé orienté
\[
\Gamma_{\infty,3}^{\circ}=\{
\underset{\substack{\Vert\\r_\infty r_1}}{\rho_s},
\underset{\substack{\Vert\\r_0r_\infty}}{\sigma},
\underset{\substack{\Vert\\r_1r_0}}{\rho_f}
\mid\rho_f\sigma\rho_s=\sigma^2=\rho_f^3=1\}.
\]
)

\margin{c'est aussi le groupe
$\widetilde{\Pi}_{0,3}/(\ell_1^2,\ell_\infty^3)$, où
$\widetilde{\Pi}_{0,3}=\Pi_{0,3}\cdot\{1,\tau\}$}

7) $\Pi_1\bigl(\mathcal U_{0,3}^{\mathrm{conf}},\operatorname{pt}(U_{0,3},-\bar j)\bigr)$
\inserted{surface conforme universelle de type $(0,3)$}
avec sous-groupe $\Pi_1\bigl(\mathcal U_{0,3}^{\mathrm{an}},\operatorname{pt}(U_{0,3},-\bar j)\bigr)$
\inserted{surface holomorphe universelle de type $(0,3)$}

\margin{\uline{NB} $\mathcal U_{0,3}^{\mathrm{conf}}\simeq(U_{0,3},\mathcal T_{0,3})$
ce qui permet d'identifier 7) à 4)}

8) $\mathcal T'_{1,1}=\mathcal T_{1,1}/\underbrace{\{\pm1\}}_{\text{centre}}$,
avec $\mathcal T_{1,1}\simeq\operatorname{Autext}_{\mathrm{lac}}(\Pi_{1,1})
\simeq A_{1,1}/A_{1,1}^{0}$

($A_{1,1}\simeq\text{autom. top. de }\unsure{U_{1,1}^{\circ}|U_{1,1}}$)
avec sous groupe $\mathcal T_{1,1}^{\prime+}=\mathcal T_{1,1}^{+}/\{\pm1\}$

\scanpage{96}
\mspage{337}

9) $\Pi_1(M_{1,1}^{\mathrm{conf}},\mathrm{pt})' \simeq \unsure{\Pi_1()/\mathrm{centre}\{\pm1\}} \simeq \Pi_1(M_{1,1}'{}^{\mathrm{conf}},\mathrm{pt})$

où $M_{1,1}^{\mathrm{an}}$ est la multiplicité modulaire \deleted{analytique}\inserted{\unsure{conf|conf.}} des courbes elliptiques \inserted{conformes} (i.e. courbes \deleted{analytiques}\inserted{conformes} de type $1,1$), $M_{1,1}'{}^{\mathrm{an}}$ celle des courbes elliptiques \deleted{analytiques}\inserted{conformes} ``mod. symétrie''; avec sous-groupe
\[
\Pi_1(M_{1,1}^{\mathrm{an}},\mathrm{pt})' \simeq \unsure{\Pi_1()/\mathrm{centre}\{\pm1\}} \simeq \Pi_1(M_{1,1}'{}^{\mathrm{an}},\mathrm{pt})
\]
où le pt base est la courbe \inserted{ell. conforme \unsure{resp.}} analytique, quotient de $\mathbb C$ par $\mathbb Z+\mathbb Z(-\bar j)$. \uline{NB} Les multiplicités modulaires sont $(D,\mathrm{Gl}(2,\mathbb Z))$, $(D,\mathrm{Sl}(2,\mathbb Z))$ dans le cadre \deleted{analytique} conforme resp. analytique pour $M_{1,1}$, et $(D,\mathrm{Gl}(2,\mathbb Z)'=\mathrm{GP}(1,\mathbb Z))$ resp. $(D,\unsure{\mathrm{Sl}(1,\mathbb Z)'|\mathrm{Sl}(2,\mathbb Z)'}=\mathrm{GP}(1,\mathbb Z)^+)$ dans le cadre conforme resp. analytique pour $M_{1,1}'$.

10) $\mathrm{Gl}(2,\unsure{\mathbb Z^{\ast}|\mathbb Z})' \simeq \mathrm{GP}(1,\mathbb Z)$, avec le sous-groupe $\mathrm{Sl}(2,\mathbb Z)'\simeq\mathrm{GP}(1,\mathbb Z)^+$.

\scanpage{97}

(VI) \uline{Les automorphismes à lacets de} $\widehat{\Pi}_{0,3}$ \uline{et de} $\widehat{\mathfrak S}_{0,3}^{+}$

On a un diagramme \unsure{comm.} canonique :

(49)
\diagram{\begin{tikzcd}[ampersand replacement=\&,column sep=large,row sep=large]
\begin{array}{l}
\operatorname{Autext}(\widehat{\mathfrak S}_{0,3}^{+})'\\
=\text{groupe des autom}\\
\text{ext. induisant}\\
\text{l'identité dans le}\\
\text{quotient }\mathfrak S_3\simeq\mathcal T_{0,3}^{+}
\end{array}
\arrow[r,"{\text{restriction}}","\sim"']
\&
\begin{array}{l}
\operatorname{Autext}^{\unsure{\rho}}(\widehat{\Pi}_{0,3})'\\
=\text{groupe des autom}\\
\text{extérieurs qui}\\
\text{commutent à}\\
\text{l'action ext. de }\mathfrak S_3
\end{array}
\\
\begin{array}{l}
\operatorname{Aut}(\widehat{\mathfrak S}_{0,3}^{+},\sigma_\infty)'\\
=\text{groupe des autom}\\
\text{de }\widehat{\mathfrak S}_{0,3}^{+}\text{ qui }\deleted{\unsure{\text{fixent}}}\inserted{\text{fixent}}\ \sigma_\infty,\\
\text{et qui fixent }\rho\ \mathrm{mod}\ \widehat{\Pi}_{0,3}
\end{array}
\arrow[u,"\sim"]\arrow[r,"{\text{restriction}}","\sim"']
\&
\begin{array}{l}
\operatorname{Aut}(\widehat{\Pi}_{0,3},\sigma_\infty)\\
=\text{groupe des aut.}\\
\text{qui commutent}\\
\text{à }\sigma_\infty,\text{ et qui}\\
\text{commutent}\\
\mathrm{ext}^{\mathrm t}.\text{ à }\rho
\end{array}
\arrow[u,"\sim"]
\end{tikzcd}}

\margin{\deleted{\illegible{1}}}

Les \unsure{$\in$} \deleted{\unsure{éléments}} $\operatorname{Aut}(\widehat{\mathfrak S}_{0,3}^{+},\sigma_\infty)'$ sont de la forme
\[
(50)\qquad
\left\{\begin{aligned}
u(\sigma_\infty)&=\sigma_\infty\\
u(\rho)&=g\rho
\end{aligned}\right.
\qquad (\text{i.e. }u\rho u^{-1}\rho^{-1}=g)
\]

\margin{On écrit les formules dans $\operatorname{Aut}(\widehat{\mathfrak S}_{0,3}^{+})'=G$, groupe des autom de $\widehat{\mathfrak S}_{0,3}^{+}$ qui induisent l'identité dans $\widehat{\mathfrak S}_{0,3}^{+}/\Pi_{0,3}\simeq\mathfrak S_3$, \uline{\unsure{s'identifie}} au groupe $\operatorname{Aut}(\Pi_{0,3})'=\mathrm{gr.}$ des autom de $\Pi_{0,3}$ qui commutent \inserted{\unsure{\uline{extérieure}}} à $\mathfrak S_3$; -- on aura $\Pi_{0,3}\subset\widehat{\mathfrak S}_{0,3}^{+}\subset\widehat{\mathfrak S}_{0,3}\subset\unsure{\mathfrak g.|g.}\,G$. \uline{NB} On admet ici que le centre de $\widehat{\mathfrak S}_{0,3}^{+}$ est trivial.}

$g\in\widehat{\Pi}_{0,3}$ \deleted{\illegible{1}} est assujetti à la condition de cocycle
\[
(51)\qquad \boxed{g\rho(g)\rho^2(g)=1}
\]
qui exprime la relation
\[
(52)\qquad (g\rho)^3=1\ .
\]

On trouve ainsi une corr. $1-1$ entre les éléments de $\operatorname{Aut}(\widehat{\mathfrak S}_{0,3}^{+},\sigma_\infty)'$ et les $g\in\widehat{\Pi}_{0,3}$ \unsure{satisfaisant} (51), et tels de plus que

(53) $\sigma_\infty,g\rho$ avec les relations $\sigma_\infty^2=(g\rho)^3=1$ \unsure{présentent} $\widehat{\mathfrak S}_{0,3}^{+}$

\scanpage{98}
\mspage{338}

(que je vais appeler la ``condition liminaire'').

\deleted{Pour \unsure{$u$ ainsi défini par un}}

Pour un $u\in G=\operatorname{Aut}(\widehat{\mathfrak S}_{0,3}^{+})'\unsure{!}$ (ne fixant pas nécessairement $\sigma_\infty$), on a des conditions équivalentes
\[
(54)\quad
\left\{\begin{aligned}
\text{(i)}\quad &u(\overset{\inserted{\varepsilon_0}}{\ell'_0})\in\text{conjugué de }\deleted{\unsure{\ell}}\varepsilon_0^{\widehat{\mathbb Z}}=L'_0\ (\text{dans }\widehat{\mathfrak S}_{0,3}^{+})\\
\text{(ii)}\quad &u(\ell_0)\in\text{conjugué de }\ell_0^{\widehat{\mathbb Z}}=L_0\ (\text{dans }\widehat{\Pi}_{0,3})\\
\text{(iii)}\quad &u(L_i)\text{ conjugué de }L_i\quad\text{pour }i=0,1,\infty\\
\text{(iv)}\quad &\exists\ p\in\widehat{\mathbb Z}^{\ast},\ g_0\in\widehat{\Pi}_{0,3}\quad
\boxed{u(\varepsilon_0)=\operatorname{int}(g_0)\varepsilon_0^p}\\
\text{(v)}\quad &\exists\ p\in\widehat{\mathbb Z}^{\ast},\ g_0\in\widehat{\Pi}_{0,3}\quad
\boxed{u(\ell_0)=\operatorname{int}(g_0)(\ell_0^p)}
\end{aligned}\right.
\]

On aura alors aussi
\[
(55)\qquad
\begin{aligned}
\deleted{\unsure{u}}\quad u(\ell_1)&=\operatorname{int}(g_1)(\ell_1^p)&
g_1&=g\rho(g_0)\\
u(\ell_\infty)&=\operatorname{int}(g_\infty)(\ell_\infty^p)&
g_\infty&=g\rho(g_1)=g\rho(g)\rho^2(g_0)\ .
\end{aligned}
\]

L'élément \deleted{$p\in\widehat{\mathbb Z}^{\ast}$ \unsure{est}}
\[
(55)\qquad p\in\widehat{\mathbb Z}^{\ast}\qquad p=\chi(u)
\]
est bien déterminé par $u$, donc par $g$\inserted{(si $u$ donné par (50))}, on l'appellera le \uline{multiplicateur} de $u$. L'élément \deleted{\unsure{$g_0$}}
\[
(56)\qquad g_0\in\widehat{\Pi}_{0,3}
\]
est déterminé a priori mod $L_0=\ell_0^{\widehat{\mathbb Z}}$ (changement
\[
(57)\qquad g_0\longmapsto g_0\ell_0^\alpha
\]
).

\margin{Notations $\operatorname{Autext}_{\mathrm{lac}}(\widehat{\mathfrak S}_{0,3}^{+})$, $\operatorname{Aut}_{\mathrm{lac}}(\widehat{\mathfrak S}_{0,3}^{+})$, $\operatorname{Aut}_{\mathrm{lac}}(\widehat{\mathfrak S}_{0,3}^{+})'$, $\operatorname{Aut}_{\mathrm{lac}}(\widehat{\mathfrak S}_{0,3}^{+},\sigma_\infty)'$, $\operatorname{Aut}_{\mathrm{lac}}(\widehat{\Pi}_{0,3},\sigma_\infty)\ldots$}

Si les conditions \inserted{équiv.} (\unsure{54|55}) sont satisfaites, on dit que $u$ respecte la structure à lacets de $\widehat{\mathfrak S}_{0,3}^{+}$ (ou de $\Pi_{0,3}$) -- en sous-entendant que $\widehat{\mathfrak S}_{0,3}^{+}$ est muni du ``groupe-lacet'' $\varepsilon_0^{\widehat{\mathbb Z}}\simeq L'_0$. On se \unsure{bornera|borne} par la suite à de tels autom. \uline{NB} Ces conditions ne dépendent que de l'autom. ext. défini par $u$, dans $\widehat{\Pi}_{0,3}$ ou dans $\widehat{\mathfrak S}_{0,3}^{+}$. On

\scanpage{99}

va \unsure{s'intéresser} aux relèvements canoniques de tels autom extérieurs donnés par (49), (50), en des autom. effectifs.

La formule de définition (50) de $u$ ne fait intervenir que $g$\inserted{$\in\widehat{\Pi}_{0,3}$}, en principe $g_0\ \mathrm{mod}\ L_0$ dans $\widehat{\Pi}_{0,3}/L_0$ est déduit de $g$. Mais en pratique $g_0$ joue un rôle aussi important que $g$ -- on verra que dans une large mesure il détermine $g$ (formule (62)).

S'il n'y avait pas lieu d'introduire $g_0$, pour exprimer les conditions (54), on pourrait aisément \uline{paramétrer} les $g$ \unsure{satisfaisant} (51) -- donc correspondants à des \deleted{\unsure{scindages de}}\inserted{relèvements de} $\dot\rho\in\mathfrak S_3$ en des éléments d'ordre $3$ de $\widehat{\mathfrak S}_{0,3}^{+}$ -- en notant qu'il y a exactement deux classes de conjugaison de tels relèvements, (correspondant aux pts fixes $-\bar j,-j$ de $\dot\rho$ opérant sur $U_{0,3}^{\mathrm{an}}$), qui se paramétrisent par
\[
(58)\quad
\left\{\begin{aligned}
g&=\lambda_0^{\unsure{\sharp|\ast}}\rho(\lambda_0^{\unsure{\sharp|\ast}})^{-1}
=[\lambda_0^{\unsure{\sharp|\ast}},\rho]
&&\begin{gathered}\text{(\deleted{première}\inserted{classe} principale, \unsure{correspondante|correspondant}}\\
\text{au ``pôle nord'' }-\bar j)\end{gathered}\\
g&=\lambda_0^{\unsure{\sharp|\ast}\prime}\ell_1^{-1}\rho(\lambda_0^{\unsure{\sharp|\ast}\prime})^{-1}
&&\begin{gathered}\text{(classe adjointe, corresp. au}\\
\text{``pôle sud'' }-j)\end{gathered}
\end{aligned}\right.
\]

\margin{\unsure{Ces|Les} \unsure{représentations} paramétriques sont \unsure{univoques|uniques}.}

où $\lambda_0,\lambda_0^{\unsure{\sharp|\ast}\prime}\in\widehat{\Pi}_{0,3}$. On verra plus bas que ces cas correspondent respectivement aux cas
\[
(59)\qquad p\equiv1\ (3)\qquad\text{et}\qquad p\equiv-1\ (3)\ .
\]

\scanpage{100}
\mspage{339}

Mais ici la situation se corse considérablement par la nécessité de tenir compte des \deleted{\unsure{la|les}} conditions (54) (p.ex. sous la forme (iv)) -- donc \inserted{celle} d'introduire conjointement $g_0,p$. Il y a lieu alors d'introduire les quantités
\[
(60)\qquad \gamma=\frac{p-1}{2}\in\widehat{\mathbb Z}
\]
(NB comme $p\in\widehat{\mathbb Z}^{\ast}$, sa \unsure{cl.} dans $\mathbb Z/2\mathbb Z$\inserted{$=\widehat{\mathbb Z}/2\widehat{\mathbb Z}$} est $1$, donc $p-1\in2\widehat{\mathbb Z}$, d'où la possibilité de trouver $\gamma$),
\[
(61)\qquad h_0=g_0^{-1}\sigma_\infty(g_0)
\]

\margin{$h_0$, relatif à l'autom. $\sigma_\infty$ d'ordre $2$, \unsure{joue} un rôle symétrique \unsure{de} $\xi$ relatif à l'autom. $\rho$ d'ordre $3$ -- ceci n'apparaît pas \unsure{non plus} dans les notations.}

\deleted{(on change nos notations, précédemment on avait pris $h_0=\sigma_\infty(g_0)^{-1}g_0$, donc on prend l'inverse \unsure{p.r.|p.v.} aux notations précédentes).} On vérifie \unsure{bien|donc} que $g$ s'exprime à l'aide de $g_0$ et de $\gamma$ (donc de $p$) à l'aide de la formule (\uline{équivalente} à (54) (iv) ou (v))
\[
(62)\quad
\boxed{g=\sigma_\infty(g_0)\ell_1^\gamma\rho(g_0)^{-1}}
\left(=(g_0\ell_0^{-\gamma})\left(\underbrace{\ell_0^\gamma h_0}_{\xi}\right)\rho(g_0\ell_0^{-\gamma})^{-1}\right)
\]

\margin{\uline{NB} Cette expression de $g$ n'est pas \unsure{changée} si $g_0\longmapsto g_0\ell_0^\alpha$.}

Donc les $u$ qui nous intéressent sont décrits par les couples $(g_0,\gamma)$ $(g_0\in\widehat{\Pi}_{0,3}/L_0,\ \gamma\in\widehat{\mathbb Z})$ \unsure{satisfaisant} les conditions
\[
(63)\qquad 2\gamma+1\in\widehat{\mathbb Z}^{\ast}\qquad (\mathrm{NB}\quad 2\gamma+1=p=\chi(u))
\]
et la condition (51), où $g$ est exprimé en fonction de $g_0$ et de $\gamma$ par (62). Cette \deleted{formule}\inserted{relation} (51) équivaut à son tour à la relation
\[
(64)\qquad\boxed{\xi\rho(\xi)\rho^2(\xi)=1}
\]
où
\[
(65)\quad\xi=\ell_0^\gamma h_0
\qquad\unsure{\deleted{\text{avec}}|\text{avec}}\ (66)\quad
\begin{aligned}
h_0&\overset{\mathrm{dfn}}{=}
g_0^{-1}\sigma_\infty(g_0)\\
&{}=[g_0^{-1},\sigma_\infty]
\end{aligned}
\]

\margin{\uline{NB} On \unsure{appelle} ici $h_0$ ce que précédemment on notait \unsure{$h_0^{-1}$}.}

\scanpage{101}

Ainsi, les $u$ sont \unsure{paramétrés} par les couples
$(g_0\in\widehat{\Pi}_{0,3}/L_0,\ \gamma\in\widehat{\mathbb Z})$, satisfaisant les conditions
(63) et (64), qui (moyennant (65)) s'écrit sous forme explicite
\[
\tag{\unsure{65bis}}
\boxed{
\bigl(\ell_0^\gamma\underbrace{g_0^{-1}\sigma_\infty(g_0)}_{h_0}\bigr)
\rho\bigl(\ell_0^\gamma\underbrace{g_0^{-1}\sigma_\infty(g_0)}_{h_0}\bigr)
\rho^2\bigl(\ell_0^\gamma\underbrace{g_0^{-1}\sigma_\infty(g_0)}_{h_0}\bigr)=1
}
\]
\margin{il vaudrait mieux considérer cette équation comme ``équation en $h_0$'' \unsure{jouant} un rôle symétrique à ``l'équation en $\xi$'' (68)}
(``équation en $g_0$'' --- par opposition avec l'``équation en $\xi$'' plus bas).
(De plus, il \unsure{faut} rajouter la condition liminaire, pour mémoire...)

On peut aussi considérer que \deleted{$g$ est} $u$ (via $g$) est défini par
$(\xi\in\widehat{\Pi}_{0,3},\ \gamma\in\widehat{\mathbb Z})$, qui permettent en effet de retrouver
\[
h_0=\ell_0^{-\gamma}\xi
\]
par (65), puis $g_0$ par (66)
\[
g_0^{-1}\sigma_\infty(g_0)=h_0
\]
qui détermine en effet $g_0$ de façon unique, grâce au fait (qu'on admet provisoirement)
que $\widehat{\Pi}_{0,3}^{\sigma_\infty}=\{1\}$.
La possibilité de trouver $g_0$ satisfaisant à cette condition équivaut (on l'admettra aussi ici) à la condition
\[
\tag{67}
\sigma_\infty(h_0)h_0=1
\]
qui, en termes de $\xi,\gamma$, s'écrit aussi sous la \unsure{forme}

\scanpage{102}
\mspage{340}

\[
\tag{68}
\boxed{\sigma_\infty(\xi)=\ell_1^\gamma\xi^{-1}\ell_0^\gamma}.
\]

Les équations (64), (68) sont appelées les ``équations en $\xi$'' --- par opposition à ``l'équation en $g_0$'' (66), qui à première vue a l'avantage d'être une seule, et le désavantage d'être plus \unsure{compliquée}.
\margin{plutôt ``l'équation en $h_0$''}
En résumé, les $u$ qui nous intéressent
$(u\in\mathrm{Aut}_{\mathrm{lac}}(\widehat{\mathfrak S}_{0,3}^{+},\sigma_\infty)')$
sont \unsure{donnés|décrits} par les couples
\deleted{$\unsure{(\xi,\gamma)}$}\inserted{$(\xi\in\widehat{\Pi}_{0,3}^{+},\ \gamma\in\widehat{\mathbb Z})$}, satisfaisant les conditions
\[
\begin{aligned}
\text{a)}\quad&2\gamma+1\in\widehat{\mathbb Z}^{*} &&(63)\\
\text{b)}\quad&\xi\rho(\xi)\rho^2(\xi)=1 &&(64)\\
\text{c)}\quad&\sigma_\infty(\xi)=\ell_1^\gamma\xi^{-1}\ell_0^\gamma &&(68)
\end{aligned}
\]
d) Condition liminaire pour mémoire (53)
(qui ne \unsure{s'expriment} apparemment pas directement de façon commode en termes de $\xi,\gamma$).

Pour que deux couples $(\xi,\unsure{\gamma})$, $(\xi',\unsure{\gamma}')$ définissent le $\overline{\mathrm m}$ $u$, il f. et s. qu'on ait
\[
\tag{69}
\unsure{\gamma}=\unsure{\gamma}',\qquad
\xi'=\ell_0^{-\alpha}\xi \ell_1^\alpha\qquad(\alpha\in\widehat{\mathbb Z}).
\]

La condition (64) peut encore se \unsure{remplacer} par une représentation paramétrique
\[
\tag{70}
\left\{
\begin{aligned}
\xi&=\lambda\rho(\lambda)^{-1}\qquad\text{ou}\\
\xi&=\lambda \ell_0^{-1}\rho(\lambda)^{-1}
\end{aligned}
\right.
\]
avec $\lambda\in\widehat{\Pi}_{0,3}$ --- auquel cas l'équation (68)

\scanpage{103}

peut s'interpréter comme une ``équation en $\lambda$'', qui suivant les deux cas prend la forme
\[
\tag{71}
\left\{
\begin{aligned}
\sigma_\infty(\lambda\rho(\lambda)^{-1})
 &=\ell_1^\gamma(\rho(\lambda)\lambda^{-1})\ell_0^\gamma\\
&\text{ou}\\
\sigma_\infty(\lambda\unsure{\ell_0^{-1}}\rho(\lambda)^{-1})
 &=\ell_1^\gamma(\rho(\lambda)\ell_0\lambda^{-1})\ell_0^\gamma,
\end{aligned}
\right.
\]

Ici pour que $(\lambda,\gamma)$ et $(\lambda',\gamma')$ définissent le $\overline{\mathrm m}$ $u$, il f et s. qu'on ait \deleted{\illegible{1}}
\[
\tag{72}
\lambda'=\ell_0^{-\alpha}\lambda,\qquad\gamma=\gamma'.
\]

\uline{Remarques} 1) Au facteur ``anodin'' $\ell_0^\gamma$ près, \unsure{qui|que} différencie $h_0$ et $\xi$ (65), on peut dire que la \unsure{recherche} des
$u\in\mathrm{Aut}_{\mathrm{lac}}(\widehat{\mathfrak S}_{0,3}^{+},\sigma_\infty)$
équivaut à celle des \deleted{\unsure{éq.}} solutions \uline{simultanées} des deux équations de cocycles
\[
\tag{73}
\left\{
\begin{aligned}
\xi\rho(\xi)\rho^2(\xi)&=1\\
h_0\sigma_\infty(h_0)&=1
\end{aligned}
\right.
\]
relatifs \deleted{respectivement} aux groupes $\mathbb Z/3\mathbb Z$, opérant sur $\widehat{\Pi}_{0,3}$ via $\rho$, et $\mathbb Z/2\mathbb Z$, opérant via $\sigma_\infty$ --- dont les solutions sous formes paramétriques sont respectivement
\[
\tag{74}
\left\{
\begin{aligned}
\xi&=\lambda\rho(\lambda)^{-1}\quad\text{ou}\quad
 \xi=\lambda \ell_0^{-1}\rho(\lambda)^{-1}\\
&\quad\text{(suivant que }\gamma\equiv0\ (3)\text{ ou }\gamma\equiv\unsure{1|-1}\ (3)\\
&\qquad\text{i.e. }p\equiv1\ (3)\text{ ou }p\equiv-1\ (3)\text{)}\\
h_0&\unsure{=|\equiv}\unsure{g_0^{-1}\sigma_\infty(g_0)}
\end{aligned}
\right.
\]
\margin{NB l'expression (62) pour $g$ met aussi bien en évidence \unsure{cette double} \unsure{intervention} de $\sigma_\infty$ et de $\rho$ dans la détermination des $u$!!!}
\margin{(NB on pourrait mettre \unsure{partout} $g_0$ pour $g_0^{-1}$...)}

\scanpage{104}
\mspage{341}

3) La relation entre le $\lambda_0$ dans la représentation paramétrique de $g$ (58), et le $\lambda$ dans celle de $\xi$ (74 a) est implicite dans la dernière expression (62) de $g$, savoir
\[
\tag{75}
\lambda_0=g_0\ell_0^{-\gamma}\lambda.
\]

4) Notons ici aussi, que les quantités
\[
\begin{array}{cccccccc}
p,&\gamma,&g,&g_0,&h_0,&\xi,&\lambda_0,&\lambda\\
\rotatebox{-90}{$\in$}&\rotatebox{-90}{$\in$}&\rotatebox{-90}{$\in$}&\rotatebox{-90}{$\in$}&\rotatebox{-90}{$\in$}&\rotatebox{-90}{$\in$}&\rotatebox{-90}{$\in$}&\rotatebox{-90}{$\in$}\\
\widehat{\mathbb Z}^{*}&\widehat{\mathbb Z}&\widehat{\Pi}_{0,3}&\widehat{\Pi}_{0,3}/L_0&&&\widehat{\Pi}_{0,3}&
\end{array}
\]
\inserted{resp. $\lambda'_0,\lambda'$}
(avec \inserted{pour $u$ constant} les indéterminations spécifiques à $g_0,h_0,\xi,\lambda$) peuvent être considérées \inserted{pour $u$ variable,} comme des fonctions sur le groupe
$\mathrm{Autext}_{\mathrm{lac}}(\widehat{\mathfrak S}_{0,3}^{+})'\simeq\mathrm{Autext}_{\mathrm{lac}}(\widehat{\Pi}_{0,3})'$
(ss-groupe des aut. ext \unsure{du groupe} à lacets $\widehat{\Pi}_{0,3}$ qui commutent à $\mathfrak S_3$), à valeurs dans les ensembles
$\widehat{\mathbb Z}^{*},\widehat{\mathbb Z},\widehat{\Pi}_{0,3},\widehat{\Pi}_{0,3}/L_0,\ ,\ ,\widehat{\Pi}_{0,3},\ $.
En particulier, on trouve ainsi 8 fonctions \inserted{\unsure{continues}} sur $\Gamma_{\mathbb Q}$, à valeurs dans les espaces appropriés --- à commencer par le caractère cyclotomique.
\margin{On verra (\illegible{2} \unsure{moins}!) que l'on peut \uline{normaliser} $g_0$, et par là aussi $h_0,\xi,\lambda,\lambda_0$, de façon à lever toute ambiguïté dans \unsure{sa|ces} \unsure{définition|définitions}...}

\scanpage{105}

Notons à ce propos les conditions équivalentes
\[
\tag{75bis}
\left\{
\begin{aligned}
\text{(i)}\quad &p=1\\
\text{(ii)}\quad &\gamma=0\\
\text{(iii)}\quad &g=\sigma_\infty(g_0)\rho(g_0)^{-1}\\
\text{(iv)}\quad &h_0=\xi\\
\text{(v)}\quad &\sigma_\infty(\xi)=\xi^{-1}
\qquad\text{(``équation en $\xi$'' simplifiée!)}\\
\text{(vi)}\quad &h_0\rho(h_0)\rho^2(h_0)=1
\end{aligned}
\right.
\]
signifiant toutes que $\chi(u)=1$. Donc la recherche des
$u\in\mathrm{Aut}_{\mathrm{lac}}(\widehat{\mathfrak S}_{0,3}^{+},\sigma_\infty)'$
\uline{de multiplicateur 1}
($u\in\mathrm{Aut}_{\mathrm{lac}}(\widehat{\mathfrak S}_{0,3}^{+},\sigma_\infty)^{\prime0}$)
équivaut donc bel et bien à celle des solutions simultanées des équations
\inserted{de cocycles} (73), ou des équations paramétriques (``de cobord'') (74).

2) Revenant au cas général (multiplicateur quelconque), le seul cas que je sache écrire explicitement
\inserted{(sauf $u=1$)} \unsure{est} celui qui correspond à
$\tau_\infty\in\unsure{\mathfrak S_{0,3}^{*}}$,
qui correspond aux caractéristiques
\[
\tag{75ter}
\left\{
\begin{aligned}
&u=\tau_\infty,\quad\chi(u)=p=-1,\quad\gamma=-1\\
&g=\ell_1^{-1}\\
&g_0=1,\quad h_0=1,\quad\xi=\ell_0^{-1}\\
&\lambda'_0=1,\quad\lambda'=\unsure{\ell_0^\gamma g_0^{-1}\lambda_0}=\ell_0^{-1}\\
&\inserted{[\,\lambda_0=\tau_\infty,\quad\lambda=\ell_0^{-1}\tau_\infty\,]}
\end{aligned}
\right.
\]
\margin{$\leftarrow$ cf. ci-dessous}

Ce serait décidément pas un luxe \unsure{d'arriver} à écrire un paquet d'autres $u$, pour se rendre compte notamment si les \unsure{foules|foule} de conditions nécessaires qu'on est amené à dégager sur les éléments provenant de $\unsure{\Gamma_{\mathbb Q}}$, ne sont pas automatiquement satisfaites!

\scanpage{106}
\mspage{342}

(VII) Calculs dans $\mathrm{Sl}(2,\widehat{\mathbb Z})'=\mathrm{Sl}(2,\widehat{\mathbb Z})/\pm1$
(\unsure{cocycles} sous $\rho$).

\margin{Calculs ``par acquit de conscience'', \unsure{qui ne serviront} peut-être plus !}

Ce groupe est un (tout petit !) quotient de $\widehat{\mathrm{Sl}(2,\mathbb Z)'}$, qui a l'avantage de se \unsure{prêter} à des calculs \inserted{algébriques} explicites commodes, en termes de matrices. Désignons par abus de langage par les $\overline{\mathrm m}$s lettres $g,g_0,h_0,\xi,\lambda_0,\lambda$ des éléments de \deleted{\illegible{1}} $\widehat{\Pi}_{0,3}$, \inserted{(ou de $\widehat{\mathfrak S}_{0,3}^{+}$ plus généralement)} et leurs images dans \inserted{$\Pi(2)'\subset$} $\mathrm{Sl}(2,\widehat{\mathbb Z})'$ \deleted{\illegible{1}} (ou dans $\mathrm{Gl}(2,\widehat{\mathbb Z})'\simeq\mathrm{Gl}(2,\widehat{\mathbb Z})/\{\pm1\}$). \deleted{Nous} \deleted{\unsure{poserons}} de façon générale, considérons une matrice
\[
(76)\qquad g=\begin{pmatrix}a&b\\c&d\end{pmatrix}\in\mathrm{Gl}(2,\widehat{\mathbb Z})
\]
on a alors
\[
(77)\quad
\left\{\begin{aligned}
\deleted{\unsure{\rho}}\rho(g)&=\rho g\rho^{-1}=\begin{pmatrix}c+d&-c\\c+d-a-b&a-c\end{pmatrix},
&\inserted{\det(g)}\rho(g^{-1})&=\begin{pmatrix}a-c&c\\a+b-c-d&c+d\end{pmatrix}\\
\rho^2(g)&=\rho^{-1}g\rho=\deleted{\unsure{(}}\begin{pmatrix}d-b&a+b-c-d\\-b&a+b\end{pmatrix},
&\inserted{\det(g)}\rho^2(g^{-1})&=\begin{pmatrix}a+b&c+d-a-b\\\deleted{\unsure{c}}b&d-b\end{pmatrix}
\end{aligned}\right.
\]
\[
(78)\quad
\left\{\begin{aligned}
\relax[g,\rho]&=g\rho g^{-1}\rho^{-1}=\begin{pmatrix}A&B\\C&D\end{pmatrix}.\ \deleted{\unsure{\frac{1}{\det(g)}}}\\
\delta A&=(a^2+b^2+ab)-(ac+bd+bc)&\delta B&=ac+bd+bc\\
\delta C&=-(c^2+d^2+cd)+(ac+bd+ad)&\delta D&=c^2+d^2+cd
\end{aligned}\right.
\]
\margin{NB $C+D-B=\unsure{1}\deleted{\unsure{\det g=\delta}}$\quad $\delta=\det g=ad-bc$ \deleted{\illegible{1}}}

Nous écrirons l'équation
\[
\deleted{(\unsure{79})}\qquad g\rho(g)\rho^2(g)=1
\]
sous la forme équivalente
\[
(79)\qquad g\rho(g)=\rho^{-1}(g^{-1})\bigl(=\rho^2(g^{-1})\bigr)
\qquad\text{\uline{dans $\mathrm{Gl}(2,\widehat{\mathbb Z})'$}}
\]
où
\[
(80)\qquad g\rho(g)=\begin{pmatrix}P&Q\\R&S\end{pmatrix}
\]

\scanpage{107}

\[
(81)\quad
\left\{\begin{aligned}
P&=(a+b)(c+d-b)&Q&=-a(c-b)-bc\overset{\inserted{\deleted{\text{si }\det g=+1}}}{=}\deleted{\unsure{1}}\delta-a(c+d-b)\\
R&=(c+d)^2-d(a+b)&S&=ad-c(c+d)=\underset{\inserted{\text{si }\det g=1}}{\delta}-c(c+d-b)\\
&\underset{\inserted{\text{si }\det g=+1}}{=}(c+d)(c+d-b)-\delta
\end{aligned}\right.
\]
où on a posé $\delta=ad-bc=\det g$.

Comparant (81) et l'expression (77) de $\rho^2(g^{-1})$, l'équation (79) -- qui s'écrit aussi
\[
g\rho(g)=\pm\rho^2(g^{-1})
\]
donne les solutions suivantes.

\margin{Quitte à remplacer $g$ par $\unsure{-g}$, \deleted{\illegible{5}} \unsure{on se ramène} à \unsure{déterminer} les solutions de l'équation}

\[
(82)\qquad \unsure{g|\rho}\rho(g)\rho^2(g)=\unsure{+|\pm}1\qquad\text{i.e.}\quad \unsure{g|\rho}\rho(g)=\unsure{+|\pm}\rho^{-1}(g^{-1})\qquad\text{\uline{dans $\mathrm{Gl}(2,\widehat{\mathbb Z})$}}
\]
\deleted{$\unsure{(82)}$ \illegible{2} $\unsure{g\rho(g)=+\rho^{-1}(g^{-1})}$.}
On distingue les solutions suivant les facteurs du produit
\[
\mathrm{Sl}(2,\widehat{\mathbb Z})'\simeq\prod_{\ell\ \mathrm{premier}}\mathrm{Sl}(2,\mathbb Z_\ell)'.
\]
\margin{\uline{NB} \unsure{Nous} \unsure{n'utiliserons|n'utilisons} cette étude de l'équation (82) que dans le cas $\det g=1$, i.e. $\delta=1$, \unsure{sauf erreurs}}

La relation
\[
\frac1\delta(a+b)=P=(a+b)(c+d-b)\qquad\text{dans }\mathbb Z_\ell
\]
donne les deux solutions $a+b=0$ ou $c+d-b=1$ (car $\mathbb Z_\ell$ est intègre).

\deleted{1°) $(a+b)\ne0$ (dans $\mathbb Z_\ell$) d'où $c+d-b=1$
\[
(83)\qquad\boxed{c+d-b=\frac1\delta}.
\]
Moyennant (83), on vérifie \unsure{aussitôt} que les formules \inserted{$g\rho(g)=\rho^{-1}(g^{-1})$} -- qu'il reste à vérifier
\[
\begin{aligned}
P&=(a+b)/\delta&Q&=(1-a)/\delta\\
R&=b/\delta&S&=(1-c)/\delta
\end{aligned}
\]
\inserted{\deleted{\illegible{1}} $\unsure{R=}$
\[
\begin{aligned}
Q&=\frac{c+d-\unsure{a}-b}{\delta}\left(=\unsure{\frac1{\delta^2}|\frac1\delta}-\frac a\delta\right),\\
R&=\frac b\delta,\quad S=\frac{d-b}{\delta}\left(=\frac1{\delta^2}-\frac c\delta\right)
\end{aligned}
\]}
(84) équivalent à \illegible{5} (82) \illegible{1}
\[
(84)\qquad\boxed{\delta^3=1}.
\]
}

\margin{La nécessité de cette formule (84) était évidente a priori.}

\deleted{2°) $a+b=0$ i.e. $a=-b$, i.e. $g=\begin{pmatrix}-b&b\\c&d\end{pmatrix}$, et
\[
-b(c+d)=1
\]
\[
(84)\quad\left\{\begin{aligned}
\det g&=-b(c+d)=\delta\quad(\text{donc }\delta+b(c+d)=0;\unsure{)}\quad d=-\frac\delta b-c\\
P&=0,&Q&=\delta+b(c+d-b)=-b^2\\
R&=(c+d)^2=\frac{\delta^2}{b^2},&S&=\unsure{ad}\ \illegible{1}\ {-c(c+d)+ad=1-c(\unsure{c+d-b})}\\
&&&\illegible{1}=\delta+c\left(\frac\delta b+b\right)
\end{aligned}\right.
\]
}

\scanpage{108}
\mspage{343}

1°) $a+b\ne0$ (dans $\mathbb Z_\ell$) d'où
\[
(83)\qquad\boxed{c+d-b=\frac1\delta}.
\]
Moyennant (83), on vérifie que les formules (82), -- \unsure{où} il reste à vérifier
\[
\begin{aligned}
Q&=\frac{c+d-(a+b)}\delta\ \unsure{\leftarrow|=}\left(=\frac1{\delta^2}-\frac a\delta\right),&R&=\frac b\delta\\
S&=\frac{d-b}\delta.\ \left(=\frac1{\delta^2}-\frac c\delta\right)
\end{aligned}
\]
équivalent à :
\[
(84)\qquad\boxed{\delta^3=1}
\]
(dont la nécessité était évidente a priori).

\margin{\uline{NB} Les conditions $c+d-b=\frac1\delta$, $\delta^3=1$ sont suffisantes \illegible{1} pour les équations (82), \unsure{!} (\unsure{indépendamment} de la question si $a+b=0$ ou $a+b\ne0$) -- elles sont nécessaires \unsure{si l'on suppose qu'on} \unsure{n'est pas} dans le cas 2° ci-dessous}

2°) $a+b=0$ (dans $\mathbb Z_\ell$) d'où
\[
(85)\qquad a=-b\qquad\text{i.e.}\quad g=\begin{pmatrix}-b&b\\c&d\end{pmatrix}
\]
\[
(86)\quad\left\{\begin{aligned}
\det g&=-b(c+d)=\delta&\text{d'où}\\
\deleted{\unsure{b+}\delta}\ \delta+b(c+d)&=0&d&=-\frac\delta b-c
\end{aligned}\right.
\]
i.e.
\[
g=\begin{pmatrix}-b&b\\c&-\frac\delta b-c\end{pmatrix},
\]
et
\[
\left\{\begin{aligned}
P&=0,&Q&=\delta+b(c+d-b)=-b^2\\
R&=(c+d)^2=\frac{\delta^2}{b^2},&S&=\delta+c\left(\frac\delta b+b\right);
\end{aligned}\right.
\]
\deleted{L'équation (82) s'écrit alors, comme \unsure{ici}}
\[
\rho^2(g^{-1})=\begin{pmatrix}0&-1/b\\b/\delta&-\frac1\delta\left(\frac\delta b+b+c\right)\end{pmatrix},
\]
\deleted{et et}

et l'équation (82) s'écrit alors
\[
(87)\quad\left\{\begin{aligned}
\text{a)}\quad&b^3=\delta^3=1.\\
\text{b)}\quad&c\bigl((b\delta)^2+(b\delta)+1\bigr)+(\delta^2+b^2\delta+b)=0.
\end{aligned}\right.
\]

\scanpage{109}

On distingue encore deux cas :

a) $b\delta=1$ i.e. $b=\frac1\delta$, alors l'équation \inserted{87 b)} devient
\[
3c+3\frac1\delta=0\qquad\text{i.e.}
\]
\deleted{$3c=$}
\[
(88)\qquad c=-\frac1\delta=\inserted{-\delta^2}\ \deleted{\unsure{3c+(2\delta^2+1)=0}}\qquad\deleted{\unsure{3c+3\delta^2=0}\quad\text{i.e.}}
\]
\deleted{Or $\delta\in\widehat{\mathbb Z}^{*}\Rightarrow\delta^2\equiv1\ (3)$ donc $2\delta^2+1\equiv0\ (3)$, d'où l'existence d'une solution dans $\widehat{\mathbb Z}$ de (88)
\[
(89)\qquad c=-\frac13(2\delta^2+1)
\]
}
correspond à
\[
(90)\qquad g=\deleted{\begin{pmatrix}
-\frac1\delta=-\delta^2&\frac1\delta=\delta^2\\
\unsure{-\frac13(2\delta^2+1)}\ \inserted{-\delta^2}&\frac13(\unsure{4|1}-\delta^2)
\end{pmatrix}}\quad
\begin{pmatrix}
-\delta'\deleted{\unsure{{}^{2}}}&\delta'\deleted{\unsure{{}^{2}}}\\
-\delta'\deleted{\unsure{{}^{2}}}&0
\end{pmatrix}
\]
\inserted{où $\delta'=1/\delta=\delta^2$, et} où, je le rappelle, $\delta$ est assujetti à la seule condition
\[
\delta^3=1.
\]

b) $b\delta\ne1$ donc (comme $(b\delta)^3=1$) $(b\delta)^2+b\delta+1=0$)
et l'équation (87 b) devient
\[
(\alpha)\qquad\delta^2+b^2\delta+b=0.\qquad\deleted{\text{\unsure{soit un}}}
\]
On aura
\[
\left\{\begin{aligned}
b&=\deleted{\unsure{\delta}}u\frac1\delta=u\delta^2\quad\text{avec}\quad u^3=1,\ u\ne1\quad\deleted{\text{donc}}\quad\text{i.e.}\\
u^2+u+1&=0
\end{aligned}\right.
\]
et l'équation $\alpha)$, \deleted{s'écrit} se réduit à
\[
\unsure{u^2|\delta^2}+u\deleted{\unsure{{}^{2}}}+1=0.
\]
\deleted{\illegible{1}}

\deleted{ce qui exige $\delta\ne1$ (sinon $u^2=0$, \deleted{\illegible{1}} incompatible avec $u^3=1$) donc $\delta^2+\delta+1=0$ -- donc $(\alpha)$ s'écrit
\[
-(\delta+1)-(u+1)+1=0\qquad\text{i.e.}
\]
\[
\delta+u+1=0
\]
ce qui \deleted{est \unsure{évidemment} \illegible{3}} \inserted{\unsure{signifie}} que $u,\delta$ sont des solutions \uline{distinctes} de l'équation cyclotomique $x^2+x+1=0$. Or $u\ne\delta$ \uline{signifie} \unsure{aussi que} $\unsure{\frac b\delta\ne\delta}$ i.e. $b\ne\delta^2$ \deleted{ce qui exige} i.e. $b=\delta$.
\[
(91)\qquad g=\begin{pmatrix}-\delta&\delta\\c&-1-c\end{pmatrix}
\]
}

\scanpage{110}
\mspage{344}

\uline{Proposition} Soit $g\in\mathrm{Gl}(2,\widehat{\mathbb Z}_\ell)$, $g=\begin{pmatrix}a&b\\c&d\end{pmatrix}$, $a,b,c,d\in\unsure{\widehat{\mathbb Z}_\ell|\mathbb Z_\ell}$,
$ad-bc=\delta\in\unsure{\widehat{\mathbb Z}_\ell^{*}|\mathbb Z_\ell^{*}}$. Pour qu'on ait
\[
g\rho(g)\rho^2(g)=1\qquad(\unsure{82|81})
\]
il f. et s. qu'on ait $\delta^3=1$, et qu'on soit dans l'un des deux cas suivants, qui s'excluent mutuellement :
\[
\begin{aligned}
1\text{°)}\quad&c+d-b=1\\
2\text{°)}\quad&g=\deleted{\left(\begin{matrix}-\delta'&\delta'\end{matrix}\right.}\ -\delta'\begin{pmatrix}-1&1\\1&0\end{pmatrix}=-\delta'\rho^{-1}\qquad\left(\text{où }\delta'=\frac1\delta=\delta^2\right)
\end{aligned}
\]

\deleted{De plus, dans le cas 1°), on a $a+b=0$ \unsure{ssi|sss} \unsure{quand de}}

\uline{Corollaire} Les $g$ pour lesquels on a \deleted{à la \uuline{fois}} \inserted{1°) et} $a+b\ne0$ sont ceux de la forme
\[
\deleted{g=\left(\begin{matrix}-\unsure{u}\delta'&u\delta'\\\unsure{c|0}&\end{matrix}\right.}
\]
\[
\left\{\begin{aligned}
g&=\begin{pmatrix}a&b\\c&b-c+\frac1\delta\end{pmatrix}\qquad\deleted{\unsure{\text{pour}}\ \unsure{c}=\frac1{(a+b)}\unsure{(a(b+\delta^2))|(ab+a\delta')}}\\
&\text{avec}\quad c=\frac1{a+b}\bigl(a(b+\delta')-\delta\bigr),
\end{aligned}\right.
\]
ceux pour lesquels on a 1°) et $a+b=0$ sont \deleted{ceux} les matrices de la forme
\[
\left\{\begin{aligned}
g&=\begin{pmatrix}-u\delta'&u\delta'\\c&-\delta'u'-c\end{pmatrix}\\
&\text{où }\delta'=1/\delta=\delta^2,\ u'=\frac1u=u^2,\quad u\text{ étant une solution}\\
&\text{de l'équation cyclotomique }u^2+u+1=0\quad\text{i.e.}\\
&\text{\unsure{racine|racines} \uline{\unsure{primitive|primitives}} }3^{\mathrm e}\text{ de l'unité}
\end{aligned}\right.
\]
